恭喜你读完了整本书！现在你可以检验一下自己的知识，试着回答几道小测验（quiz）问题并完成一些练习。

\subsection{1. 哪个 C++ 标准引入了类内默认成员初始化器？}

\begin{enumerate}
\item C++98
\item C++11
\item C++14
\item C++17
\end{enumerate}

\subsection{2. 你可以对非静态数据成员使用 auto 类型推导吗？}

\begin{enumerate}
\item 可以，自 C++11 起
\item 不可以
\item 可以，自 C++20 起
\end{enumerate}

\subsection{3. 你需要在 cpp 文件中定义静态内联（inline）数据成员吗？}

\begin{enumerate}
\item 不需要，定义就发生在声明该静态内联成员的同一位置。
\item 需要，编译器需要在 cpp 文件中提供定义。
\item 需要，编译器需要在使用该变量的所有翻译单元中都提供定义。
\end{enumerate}

\subsection{4. 静态内联变量可以是非常量吗？}

\begin{enumerate}
\item 可以，它就是一个普通变量。
\item 不可以，内联变量必须是常量。
\end{enumerate}

\subsection{5. 下面这段代码的输出是什么：}

\begin{codebox}[cpp]
struct S {
int a { 10 };
int b { 42 };
};
S s { 1 };
std::cout << s.a << ", " << s.b;
\end{codebox}

\begin{enumerate}
\item 1, 0
\item 10, 42
\item 1, 42
\end{enumerate}

\subsection{6. 考虑以下代码：}

\begin{codebox}[cpp]
class C {
C(int x) : a(x) { }
int a { 10 };
int b { 42 };
};
C c(0);
\end{codebox}

请选择正确的说法：

\begin{enumerate}
\item \texttt{C::a} 被初始化两次。第一次它用 \texttt{10} 初始化，之后在构造函数中又用 \texttt{0} 初始化。
\item \texttt{C::a} 只在构造函数中用 \texttt{0} 初始化一次。
\item 代码无法编译，因为编译器无法决定如何初始化 \texttt{C::a} 成员。
\end{enumerate}

\subsection{7. 从构造函数中抛出异常时会发生什么？}

\begin{enumerate}
\item 该对象被视为“已创建”，因此它会遵循对象的常规生命期。
\item 该对象被视为“部分创建”，因此编译器不会调用它的析构函数。
\item 编译器会调用 \texttt{std::terminate}，因为你不能从构造函数中抛出异常。
\end{enumerate}

\subsection{8. 编译这段代码时会发生什么？}

\begin{codebox}[cpp]
class Point { int x; int y; };
Point pt {.y = 10, .x = 11 };
std::cout << pt.x << ", " << pt.y;
\end{codebox}

请选择正确的说法：

\begin{enumerate}
\item 代码无法编译。指定初始化器（designated initializer）必须与 \texttt{Point} 类中数据成员的顺序一致。
\item 代码可以编译，并输出 \texttt{11, 10}。
\item 代码可以编译，并输出 \texttt{10, 11}。
\end{enumerate}

\subsection{9. 这段代码能在 C++11 中工作吗？}

\begin{codebox}[cpp]
struct User { std::string name = "unknown"; unsigned age { 0 }; };
User u { "John", 101 };
\end{codebox}

\begin{enumerate}
\item 可以，代码在 C++11 模式下能编译。
\item 代码从 C++14 模式起才能编译。
\item 即便在 C++20 中，代码也无法编译。
\end{enumerate}

\subsection{10. 下面的 \texttt{struct} 是否拥有编译器生成的拷贝构造函数？}

\begin{codebox}[cpp]
struct Test {
Test() = default;
Test(Test&& t) { /* 某个实现 */ }
int val { 10 };
};
\end{codebox}

\begin{enumerate}
\item 有，它是一个简单的类类型，因此拷贝构造函数会被隐式定义。
\item 没有，该类声明了用户定义的移动构造函数，这会阻止隐式拷贝构造函数。
\item 没有，该类声明了默认构造函数，这会阻止隐式拷贝构造函数。
\end{enumerate}

\subsection{11. 假设你有 \texttt{std::map<string, int> m;}。请从下列关于下面这个循环的说法中选出唯一正确的一项：}

\begin{codebox}[cpp]
for (const pair<string, int>& elem : m)
\end{codebox}

\begin{enumerate}
\item 该循环正确地遍历 map，不会产生额外的拷贝。
\item 由于 \texttt{elem} 的类型不匹配，该循环会为 map 中的每个元素创建一份拷贝。
\item 代码无法编译，因为 \texttt{const pair} 无法绑定到 map。
\end{enumerate}

\subsection{12. 根据 C++20，\texttt{auto x \{ 42 \};} 与 \texttt{auto z = \{ 42 \};} 是否相同？}

\begin{enumerate}
\item 相同，\texttt{x} 与 \texttt{z} 的类型都是 \texttt{int}。
\item 相同，\texttt{x} 与 \texttt{z} 的类型都是 \texttt{int \&\&}。
\item 不同，\texttt{x} 的类型是 \texttt{int}，但 \texttt{z} 的类型是 \texttt{initializer\_list<int>}。
\end{enumerate}

\subsection{13. 考虑以下代码并选出正确的说法：}

\begin{codebox}[cpp]
std::optional<std::complex<double>> opt1{std::complex<double>{0, 1}};
std::optional<std::complex<double>> opt2{std::in_place_t, 0, 1};
\end{codebox}

\begin{enumerate}
\item \texttt{opt1} 的初始化效率较低，因为我们必须创建一个临时对象，而 \texttt{opt2} 不使用任何临时对象。
\item \texttt{opt1} 的初始化与 \texttt{opt2} 一样高效；不会创建额外的拷贝。
\item 你不能在创建 \texttt{std::optional} 时使用 \texttt{in\_place\_t}。
\end{enumerate}

\subsection{14. Mayer 的单例（singleton）在 C++03 中安全吗？请选择最匹配的说法。}

\begin{enumerate}
\item 安全，它在函数作用域中使用静态变量，因此编译器会确保它在首次使用前完成初始化。
\item 安全，C++03 同时保证了线程安全与一次性初始化。
\item 不安全，只有 C++11 才为静态局部变量引入了线程安全，因此该单例模式只有从 C++11 起才是安全的。
\end{enumerate}

\subsection{15. 下面这条语句能编译吗？}

\begin{codebox}[cpp]
std::vector<std::unique_ptr<int>> ints {
std::make_unique<int>(1), std::make_unique<int>(2)
};
\end{codebox}

\begin{enumerate}
\item 不能，它无法编译，因为我们要为不可拷贝类型（\texttt{unique\_ptr}）使用一个 \texttt{initializer\_list}，而 \texttt{initializer\_list} 需要拷贝。
\item 能，它可以编译，因为 \texttt{initializer\_list} 能处理不可拷贝类型。
\item 能，它可以编译，因为编译器可以消除那些额外的拷贝。
\end{enumerate}

请写下你的答案，并在附录 B 中核对。

\section{练习}

用四个编程练习来检验你的技能。

\subsection{练习 1：NSDMI}

下面是一个包含两个数据成员的 \texttt{Point} 类声明。

\begin{codebox}[cpp]
struct Point {
double x;
double y;
};
\end{codebox}

更新这个类，让它使用 NSDMI，并将 \texttt{Point::x} 初始化为 \texttt{1.0}、将 \texttt{Point::y} 初始化为 \texttt{2.0}。

下面是测试用例的代码：

\begin{codebox}[cpp]
TEST(PointTest, X) {
Point p;
EXPECT_DOUBLE_EQ(1.0, p.x);
}
TEST(PointTest, Y) {
Point p;
EXPECT_DOUBLE_EQ(2.0, p.y);
}
\end{codebox}

你可以用下面这个 Compiler Explorer 上的解答来练习：Point tests @Compiler Explorer\footnote{\url{https://godbolt.org/z/Gor1Y8qx7}}。

运行时你会看到测试失败：

你的任务是改进代码，让测试通过：

\myGraphic{0.9}{images/fig_p258_1.png}{}

\subsection{练习 2：NSDMI}

我们再试另一个用例。下面是一个名为 \texttt{SalesRecord} 的结构体。

\begin{codebox}[cpp]
#include <string>
constexpr unsigned int DEFAULT_CATEGORY = 4;
constexpr unsigned int DEFAULT_FLAGS = 0x0a;
struct SalesRecord {
std::string name_;
double price_;
unsigned int category_ : 4;
unsigned flags_ : 4;
};
\end{codebox}

使用 NSDMI 将数据成员初始化为以下值：

\begin{itemize}
\item \texttt{name\_} 应为 \texttt{"empty"}
\item \texttt{price\_} 应为 \texttt{1.0}
\item \texttt{category\_} 应为 \texttt{DEFAULT\_CATEGORY}
\item \texttt{flags\_} 应为 \texttt{DEFAULT\_FLAGS}
\end{itemize}

下面是待解决的测试代码：

\myGraphic{0.9}{images/fig_p259_1.png}{}

\begin{codebox}[cpp]
TEST(SalesRecord, Name) {
SalesRecord s;
EXPECT_EQ("empty", s.name_);
}
TEST(SalesRecord, Price) {
SalesRecord s;
EXPECT_DOUBLE_EQ(1.0, s.price_);
}
TEST(SalesRecord, Category) {
SalesRecord s;
EXPECT_EQ(DEFAULT_CATEGORY, s.category_);
}
TEST(SalesRecord, Flags) {
SalesRecord s;
EXPECT_EQ(DEFAULT_FLAGS, s.flags_);
}
\end{codebox}

你可以用下面这个 Compiler Explorer 上的解答来练习：Point tests @Compiler Explorer\footnote{\url{https://godbolt.org/z/Y19jMs4Gb}}。

\subsection{练习 3：内联变量}

我们可以把关于构造函数与内联变量的知识结合起来，继续处理第 11 章（非局部对象）中引入的 \texttt{CountedType}。请为其它构造函数实现支持，使下面这个测试能够通过。

\begin{codebox}[cpp]
struct CountedType {
static inline int instanceCounter = 0;
static inline int maxInstanceCounter = 0;
// 在此处编写你的代码……
};
\end{codebox}

下面是测试代码：

\begin{codebox}[cpp]
int main() {
{
CountedType c0;
CountedType c1;
Tests::Expect(2, CountedType::instanceCounter);
Tests::Expect(2, CountedType::maxInstanceCounter);
CountedType c2(c1);
CountedType c3(c1);
Tests::Expect(4, CountedType::instanceCounter);
Tests::Expect(4, CountedType::maxInstanceCounter);
CountedType c4(std::move(c1));
Tests::Expect(5, CountedType::instanceCounter);
Tests::Expect(5, CountedType::maxInstanceCounter);
}
Tests::Expect(0, CountedType::instanceCounter);
Tests::Expect(5, CountedType::maxInstanceCounter);
}
\end{codebox}

如你所见，这个示例创建了若干 \texttt{CountedType} 实例，然后（通过 \texttt{Test::Expect}）检查计数器是否正确。

从下面这个可运行的代码示例 @Wandbox\footnote{\url{https://wandbox.org/permlink/GuGzTWKF8irN2YLz}} 开始，点击「Clone \& Edit」来启动该示例并进行修改。

\myGraphic{0.9}{images/fig_p261_1.png}{}

\subsection{练习 4：修正代码}

看看下面的代码，修正 \texttt{SalesRec}、\texttt{addPromo} 和 \texttt{computeTotal} 中导致测试失败的问题。

\begin{codebox}[cpp]
struct SalesRec { std::string name_;
double price_; };
void addPromo(std::vector<SalesRec>& sales, double discount) {
for (auto elem : sales)
elem.price_ = (1.0-discount)*elem.price_;
}
double computeTotal(const std::vector<SalesRec>& sales) {
double sum;
for (auto elem : sales)
sum += elem.price_;
return sum;
}
TEST(computeTotal, empty) { ... }
// 失败……
TEST(computeTotal, several) { ... } // 失败……
TEST(addPromo, simple) { ... }
// 失败……
TEST(addPromo, two) { ... }
// 失败……
\end{codebox}

这是起始的代码示例 @Compiler Explorer\footnote{\url{https://godbolt.org/z/54dh8eh3h}}。
